\documentclass[10pt,a4paper]{amsart}
\usepackage[margin=19mm]{geometry}
\usepackage{amsmath,amssymb,mathtools}
\usepackage[scr=rsfs,cal=euler]{mathalfa}
\usepackage{xfrac}
\usepackage[unicode,hidelinks,bookmarks=false]{hyperref}
\newcommand{\ie}{i.e.\ }
\newcommand{\cf}{cf.\ }
\newcommand{\resp}{respectively\ }
\newcommand{\Subsection}{\subsection}
\providecommand{\nobreakdash}{\nobreak\hbox{-}}
\setlength{\emergencystretch}{2em}
\multlinegap0pt
\usepackage[shortlabels]{enumitem}
\usepackage{tikz}
\usepackage[normalem]{ulem}
\usepackage{amssymb}
\usepackage{xcolor}
\newtheorem{proposition}{Proposition}
\newtheorem{assumption}{Assumption}
\newtheorem{lemma}{Lemma}
\newtheorem{theorem}{Theorem}
\newcommand{\CR}{\mathcal R}
\newcommand{\ac}[1]{\begin{quotation}\textbf{Sasha's comment:\
}{\textit{#1}}\end{quotation}}
\newcommand{\be}{\begin{equation}}
\newcommand{\br}{\mathbb{R}}
\newcommand{\dd}{\mathrm{d}}
\newcommand{\de}{\delta}
\newcommand{\ee}{\end{equation}}
\newcommand{\ga}{\gamma}  
\newcommand{\s}{\mathcal S}
\newcommand{\vv}{\mathbf{v}}
\newcommand{\vx}{\mathbf{x}}
\newcommand{\vy}{\mathbf{y}}
\title{Microlocal analysis of non-linear artifacts in cone beam CT}
\author{James W. Webber, Alexander Katsevich}
\date{Source: arXiv:2608.19233v1 (CC BY 4.0)}
\begin{document}
\maketitle

\noindent\emph{Source and licence.} Microlocal analysis of non-linear artifacts in cone beam CT. Version arXiv:2608.19233v1 (CC BY 4.0), distributed by arXiv under the Creative Commons Attribution 4.0 International licence, http://creativecommons.org/licenses/by/4.0/, as recorded in the arXiv licence metadata for this exact version and shown on the abstract page https://arxiv.org/abs/2608.19233v1. Full licence notice: \texttt{licenses/arxiv-2608.19233-CC-BY-4.0.txt}.\par\smallskip

\noindent\textit{Editorial scope.} Counted unit: the article's theorem log\_sing (source0.tex lines 905--917). The article's complete original proof of it is delivered with it (source0.tex lines 918--959, one \texttt{proof} environment of 42 lines). No mathematics is written, repaired or re-derived: the statement and the proof are the article's own lines, in the article's own notation, and no retained line is substituted or re-flowed.  Retained uncounted as the source-local context the proof needs: lines 623--630; lines 822--829; lines 876--882.  Nothing was omitted from any retained span, so the package carries no omission record.  No retained line contains a citation, so the wrapper carries no bibliography and no bibliography entry.  No retained line contains a figure environment, an image-inclusion command or drawing code, so no image is delivered and the package declares no asset dependency.  Editorial formatting only, in addition: an AMS \texttt{amsart} preamble that restates the article's own declarations and macros used by the retained lines, namely the environments \texttt{proposition}, \texttt{assumption}, \texttt{lemma}, \texttt{theorem} on separate counters, together with the standard packages those lines need.  The retained environments print in this document as Proposition 1, Assumption 1, Lemma 1, Theorem 1 in order of appearance, whereas the article prints the same items with the numbering of its own full document.  The complete native source of the article as distributed in the arXiv e-print for this exact version is bundled unchanged as \texttt{sources/208070/source0.tex}.
\begin{proposition}
\label{prop_sqrt}
Let $s - \phi(\vv) = 0$ be a local defining equation of a smooth two-dimensional manifold $\Gamma\subset\br^3$, where $\phi$ is smooth, $s\in \br$ and $\vv \in \br^2$. For any $r >0$ and any smooth function $u$, one has
\be
\label{dist_order_0}
u(s,\vv)(s - \phi(\vv))^r_\pm \in \dot{I}^{-(r+1)}(\Gamma)=I^{-r-5/4}(N^*\Gamma).
\ee
\end{proposition}

\begin{assumption}\label{ass:sngl}$ $
Suppose the line $L_{\vy_0}$ is tangent to $\s$ at $\vx_0$ where $\s$ is smooth.
\begin{enumerate}[label=(\arabic*)]
%\item\label{sng tang} The line $L_{\vy_0}$, where $\vy_0=(s_0,\Theta_0)\in Y$, is tangent to $\s$ only at $\vx_0$ where $\s$ is smooth.
\item\label{sng curv} The normal curvature of $\s$ along $L_{\vy_0}$ is not zero. 
\item\label{sng transv} The plane $\Pi(\vx_0,\xi)$ is not tangent to $\ga$ at $y(s_0)$, where $\xi\in N_{\vx_0}^*\s$.
\end{enumerate}
\end{assumption}

\begin{lemma}\label{lem:sing-ty} Suppose Assumption~\ref{ass:sngl} holds and $\vx_0$ is the only point where $L_{\vy_0}$ is tangent to $\s$. One has 
\be
\label{dist_R}
R\mu_E(\vy) =  g_1(\vy,E) + g_2(\vy,E)(s - \phi(\Theta))^{1/2}_\imath ,\ \vy\in B_{\vy_0}(\delta),
\ee
where $g_1\geq 0$ and $g_2$ are smooth; $g_2(\vy,E) \neq 0$ for $\vy\in\Gamma$; $\imath = +$ or $-$ depending on the orientation of $\s$ relative to $L_{\vy_0}$; $s = \phi(\Theta)$ is the local equation of $\Gamma$ near $\vy_0$, and $\phi$ is smooth.  
\end{lemma}

\begin{theorem}
\label{log_sing}
Let $R\mu_E$ take the form \eqref{dist_R} locally near $\vy_0$. Then, for $\delta$ small enough, we have
\be
\label{sqrt_sing}
\mathcal{R}\mu(\vy) = h_0(\vy) + h_1(\vy) (s - \phi(\Theta))^{1/2}_\imath + h_2(\vy) (s - \phi(\Theta))_\imath,\ \vy\in B_{\vy_0}(\delta),
\ee
where $h_i$, $i=0,1,2$, are smooth. Further,
\be\label{Rmu distr}
\CR\mu \in \dot{I}^{-3/2}(\Gamma)=I^{-7/4}(N^*\Gamma)
\ee
locally near $\vy_0$.
\end{theorem}

\begin{proof}




%In many CT applications, it is typical to take the negative log of the data before applying a reconstruction algorithm. This is because, when $\mu$ is independent of $E$, taking the log removes the non-linearity caused by the exponential. \ac{{Move this paragraph higher to where the logarithm is taken for the first time.}}


Introduce the function
\be
H(\vy,p) = -\log\Big(\int_0^{E_{\mathrm{mx}}}P(E) e^{-(g_1(\vy,E) + g_2(\vy,E)p)} \dd E\Big),\ \vy\in B_{\vy_0}(\delta),p\in\br.
\ee
Then
\be
\label{nlR 1}
\CR\mu(\vy) = H(\vy,(s - \phi(\Theta))^{1/2}_\imath),\ \vy\in B_{\vy_0}(\delta).
\ee
Since, {by our physical assumptions,} $P(E) \geq 0$ is integrable {and $\int_0^{E_{\text{mx}}}P(E)\mathrm{d}E > 0$,} it is clear that $H$ is a smooth function of $\vy$ and $p$.  

Consider the even and odd parts of $H(p)$:
\be\label{tH 1st}
H_e(p):=[H(p)+H(-p)]/2,\quad
H_o(p):=[H(p)-H(-p)]/2,
\ee
where the dependence of $H$ on $\vy$ is temporarily suppressed for simplicity. Introduce
\be\label{tq defs 1st}
h_1(p):=H_o(p)/p,\quad h_2(p):=(H_e(p)-H(0))/p^2.
\ee
Clearly, $h_1(p)$ and $h_2(p)$ are smooth and even. The $h_i$ are still smooth at zero since the leading order terms of the Taylor expansion of $H_o$ and $H_e$ are $p$ and $p^2$, respectively. Let $h_i(\vy)$, $i=1,2$, be the functions obtained by substituting $p=|s - \phi(\Theta)|$ into \eqref{tq defs 1st}. By construction, $h_i(\vy)$ are smooth in $B_{\vy_0}(\de)$. Moreover, \eqref{nlR 1}--\eqref{tq defs 1st} imply
\be
\label{nlR 2}
\CR\mu(\vy) = H(\vy,0)+h_1(\vy)(s - \phi(\Theta))^{1/2}_\imath + h_2(\vy)(s - \phi(\Theta))_\imath,\ \vy\in B_{\vy_0}(\delta),
\ee
which proves \eqref{sqrt_sing} with $h_0(\vy)=H(\vy,0)$.

Now, by Proposition \ref{prop_sqrt}, we have 
\be
\label{dist_order}
g(\vy) (s - \phi(\Theta))^r_\imath \in \dot{I}^{-(r+1)}(\Gamma)=I^{-(r+5/4)}(N^*\Gamma)
\ee
locally near $\vy_0$, for any smooth $g$. The second claim of the theorem follows from \eqref{sqrt_sing} after setting $r = 1/2$ in \eqref{dist_order}.
\end{proof}

\end{document}
